\chapter{Où vivent les plateformes crypto}\label{ch:nw:where-the-crypto-venues-live}

De nombreuses plateformes crypto parmi les plus grandes font tourner leurs moteurs d'appariement dans une région cloud publique ou à proximité immédiate, plutôt que dans leur propre
centre de données, et pour plusieurs d'entre elles le fournisseur de cloud le dit lui-même. Un billet de blog d'AWS daté de juin 2026 place Binance à Tokyo, Bybit à Singapour, Coinbase aux États-Unis et Deribit à Londres~;
une étude rapportée par CoinDesk en mars 2026 situe les validateurs d'Hyperliquid dans la région Tokyo d'AWS, aux côtés de l'infrastructure de Binance et de KuCoin~;
BitMEX a déplacé sa plateforme de la région irlandaise d'AWS vers Tokyo, parce que plus de 200~millisecondes la séparaient des sociétés qui la négociaient. Une
société qui veut être proche lance des machines dans la région de la plateforme, les mesure, et conserve les quelques-unes qui se trouvent être les plus proches
des serveurs de la plateforme.

Le chapitre~16 a décrit le cloud~; le présent chapitre porte sur la manière d'y trouver une plateforme. Ce que l'on sait publiquement de l'emplacement des plateformes~; comment les repérer
lorsqu'elles ne sont pas publiées (les plages d'adresses du fournisseur, les allers-retours depuis des lieux connus)~; la loterie du lancement de machines, jusqu'à
ce qu'une des machines tombe à proximité de la plateforme~; et la nouvelle mesure qui repère le moment où la plateforme déménage. Tout ce qui n'est pas une déclaration de la plateforme ou du fournisseur lui-même est une
simulation étiquetée.

\section{Les emplacements publiquement connus}

\begin{dated}{2026-09}{Où se trouvent les plateformes crypto hébergées dans le cloud, d'après des sources publiques}\label{dat:nw:where-the-crypto-venues-live:venues}
\begin{itemize}
\item AWS (juin 2026)~: «~Major exchanges operate across Regions: Binance in Tokyo, Bybit in Singapore, Coinbase in the US, and Deribit in London.~»
\item Recherche de Glassnode rapportée par CoinDesk (mars 2026)~: les 24~validateurs d'Hyperliquid répartis sur plusieurs \omterm{def:nw:trading-in-a-public-cloud:azid}{zones de disponibilité}
  de la région ap-northeast-1 (Tokyo) d'AWS, son API derrière AWS CloudFront~; Hyperliquid, Binance et KuCoin concentrent leur infrastructure dans ap-northeast-1~;
  les utilisateurs de Tokyo atteignent les validateurs en 2 à 3~millisecondes, les utilisateurs européens en plus de 200.
\item Étude de cas d'AWS~: BitMEX a déplacé sa plateforme de trading de la région Europe (Irlande) vers la région Asie-Pacifique (Tokyo), avec une bascule de dix minutes,
  l'exécution des ordres en dessous de 10~millisecondes et une liquidité en hausse de plus de 185~\% sur les principaux contrats.
\end{itemize}
\end{dated}

La boîte est courte parce que les sources sont peu nombreuses. Les plateformes publient rarement leur région dans leur propre documentation~; les fournisseurs les mentionnent
dans des billets de blog et des études de cas~; les chercheurs les déduisent à partir de mesures. Une région n'est d'ailleurs qu'un point de départ~: le chapitre~16 a montré qu'à l'intérieur de celle-ci la zone
compte pour des centaines de microsecondes, et l'identifiant de zone de la plateforme est publié plus rarement encore.

\begin{definition}[Système autonome, anycast, réseau de périphérie]\label{def:nw:where-the-crypto-venues-live:as}
Un \emph{système autonome}\index{système autonome} (AS) est un groupe de préfixes IP exploités par un ou plusieurs opérateurs de réseau sous une seule politique de routage,
identifié par un numéro dans le système de routage d'Internet. L'\emph{anycast}\index{anycast} est la pratique qui consiste à rendre une même adresse de service
disponible à plusieurs endroits à la fois, de sorte que chaque paquet atteint celui qui est le plus proche dans le système de routage. Un \emph{réseau de
périphérie}\index{réseau de périphérie} (edge) est l'ensemble des \omterm{def:nw:the-european-map:pop}{points de présence} d'un fournisseur situés près des utilisateurs, souvent atteints par anycast, qui terminent les connexions
et transmettent les requêtes aux serveurs d'origine.
\end{definition}

L'\omterm{def:nw:where-the-crypto-venues-live:as}{anycast} et les réseaux de périphérie expliquent pourquoi la consultation de l'adresse publique d'une plateforme ne dit souvent rien de son moteur d'appariement. L'API d'Hyperliquid, dans le
rapport ci-dessus, est atteinte par le \omterm{def:nw:where-the-crypto-venues-live:as}{réseau de périphérie} du fournisseur, alors que ses validateurs se trouvent dans une seule région~; un utilisateur en Europe résout le même nom
qu'un utilisateur à Tokyo, atteint une périphérie proche, et attend tout de même l'aller-retour vers Tokyo qui se trouve derrière.

\section{Comment découvrir un emplacement}

Un fournisseur qui publie ses plages d'adresses fournit le premier indice. AWS publie un fichier, ip-ranges.json, qui liste chaque préfixe avec sa région et
le service qui l'utilise. Rapprocher les adresses d'extrémité d'une plateforme de ce fichier, par plus long préfixe correspondant, indique à quelle région appartient une adresse, ou
qu'elle appartient au \omterm{def:nw:where-the-crypto-venues-live:as}{réseau de périphérie}, auquel cas elle ne dit rien de l'origine.

\omcode{code/firm/venuefind/firm_venuefind.py}{38}{53}{Un fichier de plages IP publié, chargé et rapproché par plus long préfixe.}\label{lst:nw:where-the-crypto-venues-live:match}

\begin{table}[ht]
\centering\small
\begin{tabular}{llll}
\toprule
Adresse d'extrémité & Plus long préfixe correspondant & Région & Service \\
\midrule
192.0.2.200 & 192.0.2.128/25 & ap-northeast-1 & EC2 \\
198.51.100.7 & 198.51.100.0/24 & ap-southeast-1 & EC2 \\
203.0.113.9 & 203.0.113.0/24 & GLOBAL & CLOUDFRONT \\
2001~:db8:1abc~::1 & 2001~:db8:1000~::/36 & ap-northeast-1 & EC2 \\
\bottomrule
\end{tabular}
\caption{Quatre adresses d'extrémité rapprochées de l'exemple du chapitre, un fichier de plages IP publié (au format documenté~; adresses tirées des
plages réservées à la documentation, aucun hôte réel n'est donc nommé). La troisième adresse appartient au \omterm{def:nw:where-the-crypto-venues-live:as}{réseau de périphérie}~: la recherche
localise la périphérie, et non la plateforme. Données~: \texttt{nw\_find.lookups()}.}\label{tab:nw:where-the-crypto-venues-live:lookups}
\end{table}

\begin{definition}[Triangulation par aller-retour]\label{def:nw:where-the-crypto-venues-live:tri}
La \emph{triangulation par aller-retour}\index{triangulation par aller-retour} estime où se trouve un serveur à partir des allers-retours mesurés vers lui depuis des sondes placées à des endroits connus~:
chaque aller-retour borne la distance (la lumière ne peut parcourir plus de la moitié de celle-ci dans la fibre) et, divisé par un \omterm{def:nw:the-north-american-map:factor}{facteur de tracé}, l'estime~;
l'estimation est le point dont les distances concordent le mieux avec ces estimations.
\end{definition}

\begin{proposition}[La borne du temps de lumière]\label{prop:nw:where-the-crypto-venues-live:bound}
Un serveur qui répond à une sonde en un aller-retour $T$, dont $E$ est le temps d'équipement et de traitement sur le trajet, se trouve à moins de $(T/2 - E)\,c_0/n_g$ de
la sonde le long de la géodésique, pour le plus petit \omterm{def:nw:the-north-american-map:floor}{indice de groupe} $n_g$ sur le chemin.
\end{proposition}

\begin{proof}
La moitié de l'aller-retour, diminuée du temps d'équipement, est le temps maximal que le signal a pu passer à voyager dans un sens~; à la vitesse $c_0/n_g$ il couvre au plus cette
distance, et aucun chemin n'est plus court que la géodésique.
\end{proof}

La borne est exacte~; l'estimation ne l'est pas. Dans l'exemple du chapitre, des sondes à Hong Kong, Singapour, Sydney et Mumbai mesurent des allers-retours de
36.5, 67.4, 98.8 et \qty{85.5}{\milli\second} vers un serveur caché à Tokyo (simulation étiquetée, avec des chemins 30~\% plus longs que les géodésiques
et un peu de bruit). La triangulation avec le bon \omterm{def:nw:the-north-american-map:factor}{facteur de tracé}, 1.3, place le serveur à moins de \qty{23}{\kilo\metre} de la vérité~; avec 1.2 ou 1.4, elle
se trompe de plus de \qty{540}{\kilo\metre}. La triangulation depuis d'autres villes trouve la région~; à l'intérieur de la région, la société doit
mesurer depuis l'intérieur.

\begin{omfigure}
\begin{tikzpicture}
\begin{axis}[axis equal image,width=11cm,xmin=-6200,xmax=6600,ymin=-6200,ymax=3200,xlabel={km à l'est de 120\textdegree{} E},
  ylabel={km au nord de 15\textdegree{} N},tick label style={font=\footnotesize,/pgf/number format/1000 sep={\,}},label style={font=\small},grid=major,grid style={gray!20},
  xtick distance=2000,table/col sep=comma,scaled ticks=false]
\addplot[only marks,mark=*,mark size=2.2pt,omDef,nodes near coords,point meta=explicit symbolic,
  every node near coord/.append style={font=\scriptsize,anchor=west}] table[x=x,y=y,meta=name,restrict expr to domain={\coordindex}{0:3},col sep=comma,row sep=\\]{x,y,name,kind\\ -615.8,805.5,sonde de Hong Kong,probe\\ -1739.8,-1508.5,sonde de Singapour,probe\\ 3353.2,-5397.7,sonde de Sydney,probe\\ -5069.4,449.0,sonde de Mumbai,probe\\ 2125.1,2286.3,serveur (caché),server (hidden)\\ 2139.8,2266.8,estimation,estimate\\};
\addplot[only marks,mark=square*,mark size=2.4pt,omThm] table[x=x,y=y,restrict expr to domain={\coordindex}{4:4}]{figdata/networks/17-where-the-crypto-venues-live/triangulation.csv};
\addplot[only marks,mark=x,mark size=4pt,very thick,omMeth] table[x=x,y=y,restrict expr to domain={\coordindex}{5:5}]{figdata/networks/17-where-the-crypto-venues-live/triangulation.csv};
\node[font=\scriptsize,anchor=south east,omThm] at (axis cs:2100,2450) {serveur (caché) et estimation};
\end{axis}
\end{tikzpicture}

\omcaption{Simulation~: \omterm{def:nw:where-the-crypto-venues-live:tri}{triangulation par aller-retour} depuis quatre sondes situées dans d'autres villes vers un serveur à Tokyo, avec le bon \omterm{def:nw:the-north-american-map:factor}{facteur de tracé}~; l'estimation
(croix) tombe à \qty{23}{\kilo\metre} du serveur caché (carré). Projection locale destinée au seul dessin. Données~: \texttt{fig\_find.py},
\texttt{nw\_find.triangulate\_example()}.}\label{fig:nw:where-the-crypto-venues-live:tri}
\end{omfigure}

\omcode{code/firm/venuefind/firm_venuefind.py}{60}{81}{Triangulation par aller-retour~: une recherche sur grille du point dont les distances concordent le mieux avec les allers-retours,
sans jamais dépasser la borne de temps de lumière d'une sonde.}\label{lst:nw:where-the-crypto-venues-live:tri}

\section{La loterie des instances}

\begin{definition}[Loterie des instances]\label{def:nw:where-the-crypto-venues-live:lottery}
Une \emph{loterie des instances}\index{loterie des instances} consiste à lancer de nombreuses machines cloud dans la région et la zone de la plateforme, à mesurer
l'aller-retour de chacune vers celle-ci, à garder les plus rapides et à libérer les autres, car le fournisseur place chaque machine sans dire au
client où.
\end{definition}

Au sein d'une zone, l'emplacement d'une nouvelle machine détermine sa distance aux serveurs de la plateforme~: le même rack, la même rangée, une autre salle. Le
fournisseur ne le précise pas, si bien que la société tire au sort. Le blog d'AWS de juin 2026 indique que les \omterm{def:nw:trading-in-a-public-cloud:pg}{groupes de placement}, le réglage des adaptateurs et le contournement du noyau (kernel bypass) «~peuvent atteindre
des latences intra-région inférieures à 150~µs~», et la mesure du chapitre~16 place un aller-retour typique au sein d'une zone à \qty{270}{\micro\second}. Le modèle
prend ces deux valeurs comme médiane et comme meilleur 1~\% des lancements, avec des allers-retours lognormaux entre les deux~: une hypothèse, signalée comme telle.

\begin{proposition}[Meilleur parmi $n$]\label{prop:nw:where-the-crypto-venues-live:bestn}
Si les allers-retours des lancements sont des tirages indépendants d'une loi $F$, l'espérance du meilleur parmi $n$ vaut $\int_0^\infty (1 - F(t))^n\,dt$, quantité qui
décroît avec $n$ de plus en plus lentement. Si chaque lancement coûte $c$ et si chaque microseconde économisée vaut $v$, le nombre optimal de lancements maximise
$v\,(E[X_{(1:1)}] - E[X_{(1:n)}]) - c\,n$.
\end{proposition}

\begin{proof}
Le minimum de $n$ tirages indépendants dépasse $t$ seulement si chacun des $n$ le dépasse, ce qui a pour probabilité $(1-F(t))^n$~; l'espérance d'une variable positive
est l'intégrale de sa fonction de survie. La valeur nette est l'économie par rapport à un lancement unique, diminuée du coût des lancements.
\end{proof}

\begin{omfigure}
\begin{tikzpicture}
\begin{semilogxaxis}[width=11cm,height=5.6cm,xmin=1,xmax=160,ymin=120,ymax=290,xlabel={instances lancées (log)},
  ylabel={espérance du meilleur aller-retour (\si{\micro\second})},tick label style={font=\footnotesize},label style={font=\small},grid=major,
  grid style={gray!25},table/col sep=comma]
\addplot[omDef,thick,mark=*,mark size=1.3pt] table[x=n,y=best_us]{figdata/networks/17-where-the-crypto-venues-live/lottery.csv};
\end{semilogxaxis}
\end{tikzpicture}

\omcaption{Simulation~: espérance du meilleur aller-retour vers la plateforme parmi $n$ lancements, pour des lancements lognormaux de médiane égale à
\qty{270}{\micro\second} et un meilleur 1~\% égal à \qty{150}{\micro\second} (valeurs publiées du chapitre~16 et du blog d'AWS, utilisées comme hypothèses).
Les dix premiers lancements font gagner \qty{94}{\micro\second}~; les 150 suivants en ajoutent 47 de plus. Données~: \texttt{fig\_find.py}, \texttt{nw\_find.lottery\_curve()}.}\label{fig:nw:where-the-crypto-venues-live:lottery}
\end{omfigure}

La courbe est celle des statistiques d'ordre de la \cref{prop:nw:where-the-crypto-venues-live:bestn}~: un lancement donne en moyenne \qty{279}{\micro\second}, dix
185, quarante 157, cent soixante \qty{138}{\micro\second}. Savoir si le lancement suivant en vaut la peine dépend de la valeur
d'une microseconde. Dans l'hypothèse de 100~USD par mois et par microseconde, et un lancement qui coûte une journée d'une instance à 16 vCPU (17.14~USD au tarif du chapitre~16), le
nombre optimal est de 77 lancements, valant environ 11\,700~USD par mois nets du coût des lancements (\cref{fig:nw:where-the-crypto-venues-live:net}).

\omcode{code/firm/venuefind/firm_venuefind.py}{90}{108}{La loterie des instances~: l'espérance du meilleur parmi $n$ lancements par simulation, et le nombre de
lancements qui maximise la valeur économisée nette de leur coût.}\label{lst:nw:where-the-crypto-venues-live:lottery}

\begin{omfigure}
\begin{tikzpicture}
\begin{axis}[width=11cm,height=5.4cm,xmin=0,xmax=150,ymin=-0.5,ymax=13,xlabel={instances lancées},
  ylabel={valeur nette (milliers de USD par mois)},tick label style={font=\footnotesize},label style={font=\small},grid=major,
  grid style={gray!25},table/col sep=comma]
\addplot[omThm,thick,mark=*,mark size=1.3pt] table[x=n,y=net_k]{figdata/networks/17-where-the-crypto-venues-live/net.csv};
\draw[omMeth,dashed] (axis cs:77,-0.5) -- (axis cs:77,13) node[pos=0.1,right,font=\scriptsize]{optimum~: 77};
\end{axis}
\end{tikzpicture}

\omcaption{Simulation~: valeur mensuelle de l'aller-retour économisé par rapport à un lancement unique, nette du coût des lancements, pour une valeur supposée de 100~USD
par mois et par microseconde et 17.14~USD par lancement~; l'optimum est plat, et tout nombre de 40 à 120 lancements reste à quelques pour cent de lui. Données~:
\texttt{fig\_find.py}.}\label{fig:nw:where-the-crypto-venues-live:net}
\end{omfigure}

\section{Remesure continue}

Les gagnants de la loterie ne le restent pas. La plateforme peut se déplacer au sein de la région ou vers une autre région, comme l'a fait BitMEX~; le fournisseur peut migrer
des machines~; la maintenance peut réorganiser une zone. Une société qui n'a mesuré qu'une fois ne s'en apercevra que lorsque ses exécutions se dégraderont. Le contrôle du chapitre est le test CUSUM du livre~7
sur la médiane quotidienne de l'aller-retour de chaque machine conservée~: la somme cumulée des résidus récursifs standardisés, signalée lorsqu'elle sort
de sa bande à 5~\%.

\begin{omfigure}
\begin{tikzpicture}
\begin{axis}[width=11cm,height=5.4cm,xmin=0,xmax=91,ymin=140,ymax=205,xlabel={jour},ylabel={médiane quotidienne de l'aller-retour (\si{\micro\second})},
  tick label style={font=\footnotesize},label style={font=\small},grid=major,grid style={gray!25},table/col sep=comma]
\addplot[omDef,thick,no markers] table[x=day,y=median_us]{figdata/networks/17-where-the-crypto-venues-live/drift.csv};
\draw[gray,dashed] (axis cs:61,140) -- (axis cs:61,205) node[pos=0.92,left,font=\scriptsize]{déplacement};
\draw[omMeth,thick] (axis cs:73,140) -- (axis cs:73,205) node[pos=0.92,right,font=\scriptsize]{alerte};
\end{axis}
\end{tikzpicture}

\omcaption{Simulation~: la médiane quotidienne de l'aller-retour d'une machine vers une plateforme passe de 160 à \qty{185}{\micro\second} lorsque la plateforme se déplace le jour 61~;
le test CUSUM du livre~7 signale le changement le jour 73. Données~: \texttt{fig\_find.py}, \texttt{nw\_find.drift\_example()}.}\label{fig:nw:where-the-crypto-venues-live:drift}
\end{omfigure}

Le test CUSUM est lent par conception~: il a été conçu pour détecter un changement de moyenne sans crier au loup, et ici il lui faut douze jours pour être sûr
d'un décalage de quatre écarts-types. Une société pressée ajoute une alerte plus rapide (un saut de la médiane au-delà d'un seuil fixe, confirmé le lendemain)
et utilise le test CUSUM pour confirmer~; la loterie repart alors.

\begin{method}[Trouver et suivre une plateforme hébergée dans le cloud]\label{met:nw:where-the-crypto-venues-live:find}
\begin{enumerate}
\item Lire ce qui est public~: la documentation et les annonces de la plateforme, les blogs et études de cas du fournisseur, les travaux de recherche~; noter chaque affirmation avec sa
  source et sa date.
\item Confronter les adresses des points de terminaison de la plateforme aux plages publiées par le fournisseur~; distinguer les régions d'origine des réseaux de périphérie (edge).
\item Trianguler à partir de sondes situées dans d'autres régions pour confirmer la région~; à l'intérieur, lancer des machines dans chaque zone et mesurer.
\item Lancer la loterie dans la meilleure zone~: lancer, mesurer avec des horodatages matériels, garder les plus rapides, libérer les autres~; s'arrêter lorsque le lancement suivant vaut
  moins qu'il ne coûte.
\item Remesurer chaque jour~; signaler les changements~; repartir lorsque la plateforme ou le fournisseur se déplace.
\end{enumerate}
\end{method}

\section{Tutoriel~: trouver la plateforme, gagner la loterie}\label{tut:nw:where-the-crypto-venues-live:tut}

\begin{tutorial}
\textbf{Objectif.} Confronter des adresses, trianguler un serveur, simuler la loterie et un déplacement, sans compte cloud.
\textbf{État final~:} \cref{fig:nw:where-the-crypto-venues-live:tri,fig:nw:where-the-crypto-venues-live:lottery,fig:nw:where-the-crypto-venues-live:net,fig:nw:where-the-crypto-venues-live:drift}
et \cref{tab:nw:where-the-crypto-venues-live:lookups}.
\begin{tutsteps}
\item \textbf{Adresses.} \texttt{firm\_venuefind.load\_ranges} lit le jeu de test (format documenté, adresses de documentation)~; \texttt{match}
  (\cref{lst:nw:where-the-crypto-venues-live:match}) trouve le préfixe le plus long.
\item \textbf{Triangulation.} \texttt{nw\_find.probes()} simule des allers-retours depuis quatre sites de la table des sites~; \texttt{triangulate}
  (\cref{lst:nw:where-the-crypto-venues-live:tri}) estime le serveur.
\item \textbf{Loterie.} \texttt{expected\_best} et \texttt{best\_n} (\cref{lst:nw:where-the-crypto-venues-live:lottery}) avec \texttt{LOTTERY},
  \texttt{VALUE} et \texttt{LAUNCH}.
\item \textbf{Remesure.} \texttt{moved} encapsule \texttt{firm\_decay.cusum}.
\end{tutsteps}
\textbf{Pour aller plus loin.} Remplacer les sondes simulées par des allers-retours mesurés sur des machines que l'on exploite~; donner à la loterie une distribution
à deux niveaux (même rack ou non) et observer comment le meilleur $n$ évolue.
\end{tutorial}

\section{Construction~: le localisateur de plateformes}\label{bld:nw:where-the-crypto-venues-live:venuefind}

\begin{build}
\bfield{Rôle} Où se trouve une plateforme hébergée dans le cloud, et à quelle distance se trouvent les machines de la société~: l'entrée des chapitres~18 et~19 ainsi que des lignes cloud
du plan du chapitre~29.
\bfield{Interface} \texttt{firm\_venuefind}~: \texttt{load\_ranges}, \texttt{match}, \texttt{distance\_bound\_km}, \texttt{triangulate},
\texttt{Lottery}, \texttt{expected\_best}, \texttt{best\_n}, \texttt{moved} (\texttt{firm.decay} du livre~7 encapsulé, non modifié).
\bfield{Règles} Les affirmations publiques sur la localisation sont accompagnées de leur source et de leur date~; une adresse de périphérie n'est jamais prise pour une origine~; la triangulation respecte chaque
borne de temps de propagation de la lumière~; la distribution de la loterie est une hypothèse énoncée, jusqu'à son remplacement par des mesures.
\bfield{Tests d'acceptation} \texttt{code/firm/venuefind/tests/}~: appariement du préfixe le plus long (IPv4 et IPv6), triangulation d'un point connu et refus d'un
point impossible, monotonie de la loterie, et drapeau CUSUM avec et sans déplacement.
\bfield{Pour aller plus loin} Mesurer les allers-retours via la périphérie pour séparer son délai de celui de l'origine~; un détecteur de changement bilatéral réglé pour la rapidité.
\end{build}

\begin{omsources}
\item AWS for Industries blog, «~Ultra-low-latency cross-Region crypto trading with Avelacom and AWS~» (juin 2026)~; étude de cas AWS sur BitMEX~; la
documentation AWS de ses plages d'adresses IP.
\item CoinDesk, «~Hyperliquid traders in Tokyo get 200-millisecond edge, Glassnode research shows~» (30 mars 2026).
\item RFC 1930 (systèmes autonomes), RFC 4786 (anycast), RFC 5737 et RFC 3849 (plages d'adresses de documentation).
\end{omsources}

\section{Exercices}

\begin{exercise}[$\star$]\label{exo:nw:where-the-crypto-venues-live:1}
Une sonde mesure un aller-retour de \qty{1}{\milli\second} vers un serveur. À quelle distance au plus le serveur peut-il se trouver, dans la fibre~?
\end{exercise}

\begin{exercise}[$\star$]\label{exo:nw:where-the-crypto-venues-live:2}
Pourquoi l'appariement de l'adresse API d'une plateforme avec les plages publiées par un fournisseur ne localise-t-il parfois rien d'utile~?
\end{exercise}

\begin{exercise}[$\star$]\label{exo:nw:where-the-crypto-venues-live:3}
En utilisant la \cref{fig:nw:where-the-crypto-venues-live:lottery}, combien le dixième lancement fait-il gagner par rapport au premier, et le
cent soixantième par rapport au quarantième~?
\end{exercise}

\begin{exercise}[$\star\star$]\label{exo:nw:where-the-crypto-venues-live:4}
Pourquoi la triangulation depuis d'autres villes permet-elle de trouver une région mais pas une zone~?
\end{exercise}

\begin{exercise}[$\star\star$]\label{exo:nw:where-the-crypto-venues-live:5}
BitMEX est passé de l'Irlande à Tokyo. Qu'est-ce qui a changé pour un teneur de marché qui fait tourner ses machines en Irlande, et qu'aurait-il dû faire~?
\end{exercise}

\begin{exercise}[$\star\star$]\label{exo:nw:where-the-crypto-venues-live:6}
Pourquoi le test CUSUM est-il lent, et avec quoi l'associer~?
\end{exercise}

\begin{exercise}[$\star\star\star$]\label{exo:nw:where-the-crypto-venues-live:7}
\emph{Programmation.} Avec \texttt{firm.venuefind}, trouver le meilleur nombre de lancements si une microseconde ne vaut que 20~USD par mois.
\end{exercise}

\begin{exercise}[$\star\star\star$]\label{exo:nw:where-the-crypto-venues-live:8}
\emph{Trouver la faille.} «~Notre machine interroge l'API de la plateforme en \qty{1}{\milli\second} depuis Francfort, donc la plateforme se trouve à Francfort.~»
\end{exercise}

\section{Problème~: lancer quarante, en garder un}

\begin{problem}[{Problème du week-end --- combien de machines lancer}]\label{pb:nw:where-the-crypto-venues-live:1}
Une société négocie sur une plateforme qui fonctionne dans une zone d'une région cloud. Les allers-retours de ses lancements vers la plateforme suivent une loi lognormale de médiane
\qty{270}{\micro\second}, et leur meilleur 1~\% se situe à \qty{150}{\micro\second}. Une microseconde lui vaut 100~USD par mois~; un lancement, mesuré pendant une journée, coûte
17.14~USD.

\medskip\noindent\textbf{Partie I --- Les tirages.}
\begin{enumerate}
\item Quel est l'aller-retour espéré d'un lancement isolé~?
\item Et le meilleur espéré parmi 10, 40 et 160~?
\item Pourquoi l'aller-retour isolé espéré est-il supérieur à la médiane~?
\item Que dit la \cref{prop:nw:where-the-crypto-venues-live:bestn} sur la forme de la courbe~?
\end{enumerate}

\medskip\noindent\textbf{Partie II --- L'économie.}
\begin{enumerate}[resume]
\item Quelle est la valeur nette de la conservation du meilleur parmi 40~?
\item Quel nombre de lancements maximise la valeur nette, et que vaut-elle~?
\item Quelle est la platitude de l'optimum~?
\item Qu'en est-il si une microseconde ne vaut que 20~USD par mois~?
\end{enumerate}

\medskip\noindent\textbf{Partie III --- D'où viennent les tirages.}
\begin{enumerate}[resume]
\item Quels faits du chapitre~16 et du blog AWS le modèle utilise-t-il, et quelles hypothèses~?
\item Pourquoi les lancements ne sont-ils peut-être pas indépendants~?
\item Qu'est-ce, hormis la position de l'instance, qui modifie son aller-retour~?
\item Comment remplacer l'hypothèse par des données~?
\end{enumerate}

\medskip\noindent\textbf{Partie IV --- Le verdict.}
\begin{enumerate}[resume]
\item Énoncer le \emph{résultat nommé}~: l'aller-retour optimal espéré parmi $n$ instances et le nombre de lancements qui maximise la valeur par microseconde nette
  du coût de lancement.
\item Combien de temps le gagnant reste-t-il gagnant~?
\item Que change une plateforme qui publie l'identifiant de sa zone~?
\item La loterie est-elle équitable pour les autres utilisateurs de la plateforme~?
\item Que ferait une plateforme pour y mettre fin~?
\item Que coûte la loterie à une société qui ne participe pas à la course~?
\item Où cela s'inscrit-il dans le plan de connectivité du chapitre~29~?
\item En une phrase~: pourquoi lancer quarante et en garder un~?
\end{enumerate}
\end{problem}

\section{Questions d'entretien}

\begin{interviewq}[$\star$ \iqroles{développeur}]\label{iq:nw:where-the-crypto-venues-live:1}
Comment trouveriez-vous dans quelle région cloud une plateforme d'échange de cryptoactifs fonctionne~?
\end{interviewq}

\begin{interviewq}[$\star\star$ \iqroles{développeur}]\label{iq:nw:where-the-crypto-venues-live:2}
Qu'est-ce que l'\omterm{def:nw:where-the-crypto-venues-live:as}{anycast}, et pourquoi peut-il fausser une mesure de latence~?
\end{interviewq}

\begin{interviewq}[$\star\star$ \iqroles{développeur, chercheur}]\label{iq:nw:where-the-crypto-venues-live:3}
Vous lancez 50 instances et mesurez leurs allers-retours vers une plateforme. Combien en gardez-vous, et comment décidez-vous~?
\end{interviewq}

\begin{interviewq}[$\star\star$ \iqroles{développeur}]\label{iq:nw:where-the-crypto-venues-live:4}
Comment détectez-vous qu'une plateforme a déplacé ses serveurs~?
\end{interviewq}

\begin{interviewq}[$\star\star$ \iqroles{trader}]\label{iq:nw:where-the-crypto-venues-live:5}
Pourquoi tant de plateformes de cryptoactifs et de sociétés de trading se sont-elles retrouvées à Tokyo~?
\end{interviewq}

\begin{interviewq}[$\star\star\star$ \iqroles{développeur, chercheur}]\label{iq:nw:where-the-crypto-venues-live:6}
Estimez la position d'un serveur à partir d'allers-retours mesurés dans quatre villes. Qu'est-ce qui limite la précision~?
\end{interviewq}
